\section{Proof of Theorems \ref{2406:thm:main:sgd_weak_recovery} and \ref{2406:thm:main:no_yhat_weak_recovery}}
\label{2406:sec:app:proofs}


\subsection{Preliminaries}

\paragraph{Notations and definitions} We denote by $\polylog x$ any polynomial in $\log x$ with degree $> 1$. Since the case where $n_b = O(1)$ is already covered by the results in \cite{arous2021online}, we shall always assume that $\mu > 0$. The Hermite coefficients of $\sigma$ and $f^\star$ will be denoted by $(c_k)_{k \geq 0}$ and $(c_k^\star)_{k \geq 0}$, respectively. To break the symmetry between $m$ and $-m$ inherent to the problem, we assume without loss of generality that
\[m_0 > 0 \quad \text{and} \quad c_\ell c_\ell^\star > 0.\]

Throughout this section, the update process on $\vec{w}_t$ will be written as
\begin{equation}\label{2406:eq:app:one_step_w}
    \vec{w}_{t+1} = \frac{\vec{w}_t - \gamma  \vec{g}_t}{\norm{\vec{w}_t - \gamma  \vec{g}_t}},
\end{equation}
where $\vec{g}_t$ is the gradient at time $t$: $\vec{g}_t = \nabla_{w_{t}} \ell_t$, and $\ell_t$ is the empirical loss at time $t$, that can be either the correlation or the square loss. When considering the update of the process $(\vec{w}_t)_{t \geq 0}$, it will be useful to distinguish between the randomness in $\vec{w}_t$ and the one introduced by the batch drawn at time $t$. To this end, we introduce the filtration $(\cF_t)_{t\geq 0}$ adapted to the process $\vec{w}_t$, and we shall denote by $\dP_t$ (resp. $\dE_t$) the probability (resp. expectation) conditioned on $\cF_t$.

\paragraph{Concentration in Orlicz spaces} We first recall some fact about Orlicz spaces that will be useful for our concentration bounds.
 \begin{definition}
     For any $\alpha \in \R$, let $\psi_\alpha(x) = e^{x^\alpha} - 1$. Let $X$ be a real random variable; the \emph{Orlicz norm} $\norm{X}_{\psi_\alpha}$ is defined as
     \begin{equation}
         \norm{X}_{\psi_\alpha} = \inf \left\{t > 0 :  \dE\left[ \psi_\alpha\left(\frac{|X|}{t} \right)\right] \leq 1\right\}
     \end{equation}
 \end{definition}
 It can be checked that $\norm{\cdot}_{\psi_\alpha}$ is a well-defined norm on random variables for $\alpha \geq  1$, and can be slightly modified into a norm when $\alpha < 1$; see \cite{ledoux1991probability,vandervaart1996weak} for more information. We say that a random variable is sub-gaussian (resp. sub-exponential) if its $\psi_2$ (resp. $\psi_1$) norm is finite. The main use of this definition is the following concentration inequality:
 \begin{lemma}\label{2406:lem:app:orlicz_concentration}
    Let $X$ be a random variable with finite $\psi_\alpha$-norm for some $\alpha > 0$. Then
      \begin{equation}
     \mathbb{P}\left[\left| X - \dE X \right| > t\norm{X}_{\psi_\alpha}\right] \leq 2e^{-t^\alpha}.
    \end{equation}
 \end{lemma}
As a result, any $\psi_\alpha$-norm bound yields exponential convergence tails. Orlicz norms are also well-behaved with respect to products:
 
 \begin{lemma}\label{2406:lem:app:orlicz_submult}
     Let $X$ and $Y$ be two random variables such that $\norm{X}_{\psi_\alpha}$ and $\norm{Y}_{\psi_\beta}$ are finite for some $\alpha, \beta > 0$. Then
    \[\norm{XY}_{\psi_{\lambda}} \leq \norm{X}_{\psi_\alpha} \norm{Y}_{\psi_\beta},\]
    where $\lambda$ is the number satisfying $\frac1\alpha + \frac1\beta = \frac1\lambda$.
\end{lemma}
\begin{proof}
    Assume without loss of generality that $\norm{X}_{\psi_\alpha} = \norm{Y}_{\psi_\beta} = 1$.
    We use the following Young inequality: for any $a, b > 0$, and $p, q$ such that $\frac1p + \frac1q = 1$,
    \[ab \leq \frac{a^p}{p} + \frac{b^{q}}{q}\]
    Applying this inequality to $p= \alpha/\lambda$, $q = \beta/\lambda$, $a = X^\lambda$, $b = Y^\lambda$, we get
    \[ (XY)^\lambda \leq \frac{\lambda X^\alpha}{\alpha} + \frac{\lambda Y^\beta}{\beta}. \]
    Then
    \begin{align*}
        \exp((XY)^\lambda) &\leq \exp\left( \frac{\lambda X^\alpha}{\alpha}\right) + \exp\left(\frac{\lambda Y^\beta}{\beta} \right)\\
        &\leq \frac{\lambda}{\alpha} \exp(X^\alpha) + \frac{\lambda}{\beta} \exp(Y^\beta),
    \end{align*}
    where at the last line we used Young's inequality again with the same $p$ and $q$. The result ensues from taking expectations on both sides, and noticing that $\lambda/\alpha + \lambda/\beta = 1$ by definition.
\end{proof}

 Finally, we shall use the following theorem:
 \begin{theorem}[Theorem 6.2.3 in \cite{ledoux1991probability}] \label{2406:thm:app:orlicz_sum}
     Let $X_1, \dots, X_n$ be $n$ independent random variables with zero mean and second moment $\dE X_i^2 = \sigma_i^2$. Then,
     \begin{equation}
         \norm{\sum_{i=1}^n X_i}_{\psi_\alpha} \leq K_\alpha \log(n)^{1/\alpha} \left(\sqrt{\sum_{i=1}^n \sigma_i^2} + \max_{i}\norm{X_i}_{\psi_\alpha} \right)
     \end{equation}
 \end{theorem}

\subsection{Computing the gradient at time $t$}

Throughout this section, the update process on $\vec{w}_t$ will be written as
\begin{equation}
%\label{2406:eq:app:one_step_w}
    \vec{w}_{t+1} = \frac{\vec{w}_t - \gamma  \vec{g}_t}{\norm{\vec{w}_t - \gamma  \vec{g}_t}},
\end{equation}
where $\vec{g}_t$ is the gradient at time $t$: $\vec{g}_t = \nabla_{w_{t}} \ell_t$, and $\ell_t$ is the empirical loss $t$, that can be either the correlation or the square loss. A direct computation of both gradients implies the following lemma:
\begin{lemma}
Define
\begin{equation}\label{2406:eq:app:def_ghat_gstar}
    \vec{g}_t^\star = \frac{1}{n_b} \sum_{\nu=1}^{n_b} f^\star(\langle \vec{w}^\star, \vec{z}^\nu \rangle) \sigma'(\langle \vec{w}_t, \vec{z}^\nu \rangle) \vec{z}^\nu \quad \text{and} \quad \hat{\vec{g}}_t = \frac{1}{n_b} \sum_{\nu=1}^{n_b} \sigma(\langle \vec{w}_t, \vec{z}^\nu \rangle) \sigma'(\langle \vec{w}_t, \vec{z}^\nu \rangle) \vec{z}^\nu,
\end{equation}
    Then the gradient of the correlation loss $\ell^{\mathrm{corr}}$ is $- \vec{g}_t^\star$, while the gradient of the square loss $\ell^{\mathrm{sq}}$ is $\hat{\vec{g}}_t - \vec{g}_t^\star$.
\end{lemma}

Hence, the main difference between the gradients of the correlation and square loss is a so-called \emph{interaction term} $\hat{\vec{g}}$, that only depends on the learned vector $\vec{w}_t$. Notice that $\vec{g}_t^\star$ is an average of $n_b$ independent variables of the form
\begin{equation}
    \vec{g}_t^{\star\nu} := f^\star(\langle \vec{w}^\star, \vec{z}^\nu \rangle) \sigma'(\langle \vec{w}_t, \vec{z}^\nu \rangle) \vec{z}^\nu,
\end{equation}
and we define $\hat{\vec{g}}_t^\nu$ and $\vec{g}_t^\nu$ in the same way. By Assumption \ref{2406:assump:poly_growth} and Lemma \ref{2406:lem:app:orlicz_submult}, each variable $\vec{g}_t^{\star\nu}$ (resp. $\hat{\vec{g}}_t^\nu, \vec{g}_t^\nu$) has finite $\psi_\alpha$-norm for some $\alpha > 0$, and hence Proposition 2 of \cite{dandi2023how} holds up to $\polylog(n)$ factors. 

We can also compute the conditional expectation of the gradient $\vec{g}_t$:
\begin{lemma}\label{2406:lem:app:exp_g}
    For any $t \geq 0$, 
    \begin{equation}
        \Et{\vec{g}_t^\star} = \phi(m_t) \vec{w}_t^\star +  \psi^{\mathrm{corr}}(m_t) \vec{w}_t 
    \end{equation}
    where $\phi(m_t)$ and $\psi^{\mathrm{corr}}$ are two functions with Taylor expansion
    \begin{equation}
        \phi(m) = \sum_{k=0}^\infty c_{k+1}c_{k+1}^\star m^k \quad \text{and} \quad \psi^{\mathrm{corr}}(m) = \sum_{k=0}^\infty c_{k+2}c_{k}^\star m^k.
    \end{equation}
    Further, we have
    \begin{equation}
        \Et{\hat{\vec{g}}_t} =  c^{\mathrm{sq}} \vec{w}_t \quad \text{with} \quad c^{\mathrm{sq}} = \Ec{z \sigma(z) \sigma'(z)}{z\sim \cN(0, 1)}.
    \end{equation}
\end{lemma}

\begin{proof}
    The expectation of $\vec{g}_t^\star$ is a specialization of Lemma 4 from \cite{dandi2023how} to $r = 1$. By the independence properties of Gaussians, $\Et{\sigma(\langle \vec{w}_t, \vec{z} \rangle) \sigma'(\langle \vec{w}_t, \vec{z} \rangle) \langle\vec{z}, \vec{w}'\rangle} = 0$ as soon as $\vec{w}'$ is orthogonal to $\vec{w}_t$, hence the expectation of $\hat{\vec{g}}_t$ lies along $\vec{w}_t$, and the second result follows.
\end{proof}

In the following, we will denote $\psi^{\mathrm{sq}}(x) = \psi^{\mathrm{corr}}(x) - c^{\mathrm{sq}}$. A generic $\psi$ will be used when no specialization is necessary, so that
\begin{equation}\label{2406:eq:app:exp_g}
    \Et{\vec{g}_t} = -\phi(m_t) \vec{w}^\star - \psi(m_t) \vec{w}_t.
\end{equation}

\subsection{A differential inequality for $m_t$}

 The structure of the proof is similar to the one of \cite{arous2021online}. We define the following stopping times for $\zeta > 0$:
\begin{align}
    t^+_{\zeta} &= \min\{t \geq 0: m_t \geq \zeta\}, & \ t^-_{\zeta} &= \min\{t \geq 0: m_t \leq \zeta \},
\end{align}
and the following $\gamma$-dependent time:
\begin{equation}
    \tilde t_{\gamma, \zeta}^+ = \min\{t \geq 0: \gamma m_t^{\ell-1} \geq \zeta\}
\end{equation}
 Our first goal is to show the following high-probability inequality:
\begin{proposition}\label{2406:prop:app:diff_ineq_m}
Define
\begin{equation}
    t_{\max} = \frac{n_b}{C_{\max} d \log(d)^{C_{\max}}  \gamma^2},
\end{equation}
for some sufficiently large $C_{\max} > 0$. Then, for a sufficiently small choice of $c_\gamma$:
\begin{enumerate}
    \item If we are using the square loss, and $\gamma \leq c_\gamma (n_b d^{-\ell/2} \wedge 1)$, there exists $c, \eta > 0$ such that
    \begin{equation}
        \mathbb{P}\left( m_t \geq \frac{3}{4} m_0 + c \gamma \sum_{s = 0}^{t-1} m_s^{\ell-1} \quad \forall t \leq t^+_\eta \wedge t_{\max} \right) \geq 1 - c e^{-c\log(n)^2}.
\end{equation}
    \item If we are using the correlation loss, and $\gamma \leq c_\gamma n_b d^{-\ell/2}$, there exist $c, \varepsilon > 0$ such that
    \begin{equation}
            \mathbb{P}\left( m_t \geq \frac{3}{4} m_0 + c \gamma \sum_{s = 0}^{t-1} m_s^{\ell-1} \quad \forall t \leq t^+_\eta \wedge \tilde t_{\gamma, \varepsilon}^+ \wedge t_{\max} \right) \geq 1 - c e^{-c\log(n)^2}.
    \end{equation}
\end{enumerate}
\end{proposition}

The rest of this section is devoted to show Proposition \ref{2406:prop:app:diff_ineq_m}. We define the following ``good'' event at time $t$:
\[ \cE_t \coloneqq \left\{ |\langle \vec{g}_t, \vec{w}_t \rangle | \leq \frac1{2\gamma} \right\} \]
% We can show that $\cE_t$ is a high-probability event, under certain conditions:
% \begin{lemma}
%     If we are using the square loss and $\gamma \leq c_\gamma$ for a sufficiently small choice of $c_\gamma$, then for any $T \geq 0$
%     \[ \mathbb{P}\left[\cE_t \  \forall t \leq T\right] \geq 1 - ce^{-c\log(n)^2}. \]
%     If we are using the correlation loss, for any $T \geq 0$
%     \[ \mathbb{P}\left[\cE_t \  \forall t \leq t^+_{\varepsilon/\gamma} \wedge T \right] \geq 1 - ce^{-c\log(n)^2}. \]
% \end{lemma}

\paragraph{A (almost) deterministic update inequality}

We first expand the projection step to obtain a difference inequality for the process $m_t$. We write 
\begin{equation}\label{2406:eq:app:grad_decomp}
    \vec{g}_t = \langle \vec{g_t}, \vec{w_t} \rangle \vec{w}_t + \vec{g_t}^\bot,
\end{equation}  
where $\vec{g}_t^\bot$ is orthogonal to $\vec{w_t}$. Similarly to Equation~\ref{2406:eq:app:exp_g}, we can compute the expectation of $\vec{g}_t^\bot$:
\begin{equation}\label{2406:eq:app:exp_gbot}
    \Et{\vec{g}_t} = -\phi(m_t)(\vec{w}_t^\star - m_t \vec{w}_t)
\end{equation}
\begin{lemma}\label{2406:lem:app:difference_m_determ}
For any $t \geq 0$, there exists a (random) constant $c_t$ such that the following inequality holds:
\begin{equation}
    m_{t+1} \geq m_t - \gamma c_t \langle \vec{w}^\star,\vec{g}_t^\bot \rangle - \frac{\gamma^2 c_t^2 m_t \norm{\vec{g}_t^\bot}^2}{2} - \frac12\gamma^3 c_t^3 |\langle w^\star, \vec{g}^\bot \rangle| \norm{\vec{g}_t^\bot}^2.
\end{equation}

Further, under the event $\cE_t$, we have $1/2 \leq c_t \leq 2$.
\end{lemma}

\begin{proof}
    We use the decomposition of eq. \eqref{2406:eq:app:grad_decomp} and write
    \[ \vec{w}_t - \gamma \vec{g}_t = (1 - \gamma \langle \vec{g_t}, \vec{w_t} \rangle)\vec{w}_t - \gamma \vec{g}_t^\bot \]
    As a result, if we define
    \[ c_t \coloneqq \frac1{1 - \gamma \langle \vec{g_t}, \vec{w_t} \rangle}, \]
    we have
    \[ \vec{w}_{t+1} = \frac{\vec{w}_t - \gamma c_t \vec{g}_t^\bot}{\norm{\vec{w}_t - \gamma c_t \vec{g}_t^\bot}} \]
    since the update equation \eqref{2406:eq:app:one_step_w} is invariant w.r.t scaling. Taking the scalar product of the above with $\vec{w}^\star$, we have
    \begin{equation}\label{2406:eq:app:one_step_m}
        m_{t+1} = \frac{m_t - \gamma c_t \langle \vec{w}^\star,\vec{g}_t^\bot \rangle}{\norm{\vec{w}_t - \gamma c_t \vec{g}_t^\bot}}.
    \end{equation}
    Expanding the norm in the denominator, and using that $\norm{\vec{w}_t} = 1$ and $\langle \vec{w}_t, \vec{g}_t^\bot \rangle = 0$:
    \[ \norm{\vec{w}_t - \gamma  \vec{g}_t} = \sqrt{1 + \gamma^2 c_t^2 \norm{\vec{g}_t^\bot}^2} \]
    By the convexity inequality $(1+x)^{-1/2} \geq 1 - x/2$, valid for all $x \geq 0$, Equation \eqref{2406:eq:app:one_step_m} becomes
    \begin{align*}
        m_{t+1} &\geq \left(m_t - \gamma c_t \langle \vec{w}^\star,\vec{g}_t^\bot \rangle \right) \left( 1 - \frac{\gamma^2 c_t^2}2 \norm{\vec{g}_t^\bot}^2 \right).
    \end{align*}
    The lemma ensues upon expanding and rearranging the terms.
\end{proof}

The expansion in Lemma \ref{2406:lem:app:difference_m_determ} can be decomposed in two terms: the term linear in $\gamma$ is a noisy \emph{drift term}, that will drive the dynamics, and that we will decompose as a sum of a deterministic process and a martingale. All other terms in $\gamma^2$ or $\gamma^3$ are corrections that we bound with high probability.


\paragraph{The linear term} We first control the term linear in $\gamma$. We can write
\begin{equation}\label{2406:eq:app:martingale_decomp}
     \langle \vec{w}^\star, \vec{g}_t^\bot \rangle =  \langle \vec{w}^\star, \Et{\vec{g}_t^\bot} \rangle + Z_t,
\end{equation}
where $(Z_t)_{t \geq 0}$ is by definition a martingale difference sequence for the filtration $(\cF_t)_{t \geq 0}$. The expectation term is straightforward to compute using \eqref{2406:eq:app:exp_gbot}:
\begin{equation}
    \langle \vec{w}^\star, \Et{\vec{g}_t^\bot} \rangle = - (1 - m_t^2) \phi(m_t). \label{2406:eq:app:drift_term}
\end{equation}

The contribution of the terms $Z_t$ is bounded by the following lemma:
\begin{lemma}\label{2406:lem:app:bound_martingale}
    There exists constants $c, C > 0$ such that with probability $1-ce^{-c\log(n)^2}$,
    \begin{equation}
        \sup_{1 \leq t \leq T} \sum_{s = 1}^t Z_s  \leq \frac{C\log(d)^C \sqrt{T}}{\sqrt{n_b}}
    \end{equation}
\end{lemma}
\begin{proof}
     The martingale increment $Z_t$ is an average of $n_b$ independent terms $Z_t^\nu$, that satisfy $\norm{Z_t^\nu}_{\psi_\alpha} \leq C$ for some $\alpha, C > 0$ by Assumption \ref{2406:assump:poly_growth}. As a result, if we define
     \begin{equation*}
         B_\alpha = \sup_{t \geq 0} \norm{Z_t}_{\psi_\alpha},
     \end{equation*}
     Theorem \ref{2406:thm:app:orlicz_sum} implies that
    \begin{equation*}
        B_\alpha = \frac{\polylog(d)}{\sqrt{n_b}}.
    \end{equation*}
 We now apply Theorem F.1 in \cite{li2021stochastic} with $z = \log(d)^{\frac{2(\alpha+2)}{\alpha}}  \sqrt{T} B_\alpha$, which yields the exact bound needed. 
\end{proof}

\paragraph{Bounding the corrections} Our next step is to handle the higher-order corrections. We show the following lemma:
\begin{lemma}\label{2406:lem:app:bound_corrections}
    Let $T \geq 0$, and $\eta < 1$. There exists a constant $C > 0$ such that for any $t \leq T$
    \begin{equation}
        \mathbb{P}\left(\norm{\vec{g}_t^\bot}^2 \leq C\left(\phi(m_t)^2(1 - m_t)^2  + \frac {d\log(d)^C} {n_b}\right) \quad \forall t \leq t_\eta^+ \wedge T \right) \geq 1 - cTe^{-c\log(d)^2}
    \end{equation}
\end{lemma}


\begin{proof}
    Fix some $t \in [T]$. We can write
    \[ \vec{g}_t^\bot = a(\vec{w}^\star - m\vec{w}) + \frac1n \sum_{\nu=1}^n f^\star(\langle \vec{w}^\star, \vec{z}^\nu \rangle) \sigma'(\langle \vec{w}_t, \vec{z}^\nu \rangle) \vec{z}^{\nu\bot}\]
    where each $\vec{z}^{\nu\bot}$ is orthogonal to both $\vec{w}$ and $\vec{w}^\star$. From Lemmas 9 and 11 in \cite{dandi2023how}, with probability $1 - ce^{-c\log(d)^2}$,
    \[ \frac1n \sum_{\nu=1}^n f^\star(\langle \vec{w}^\star, \vec{z}^\nu \rangle) \sigma'(\langle \vec{w}_t, \vec{z}^\nu \rangle) \vec{z}^{\nu\bot} \leq C \frac{d\log(d)^C}{n_b}. \]
    Now, we have
    \[ \langle \vec{g}_t^\bot, \vec{w}_t^\star \rangle^2 = a^2(1 - m_t^2)^2 \quad \text{and} \quad \norm{a(\vec{w}^\star - m\vec{w})}^2 = a^2 (1-m_t^2), \]
    hence
    \[ \norm{ \vec{g}_t^\bot}^2 \leq \frac1{1 - m_t^2} \langle \vec{g}_t^\bot, \vec{w}_t^\star \rangle^2 + C \frac{d\log(d)^C}{n_b}.  \]
    It remains to notice that 
    \[ \langle \vec{g}_t^\bot, \vec{w}_t^\star \rangle^2 = (\langle \vec{w}^\star, \Et{\vec{g}_t^\bot} \rangle + Z_t)^2 \leq 2(\langle \vec{w}^\star, \Et{\vec{g}_t^\bot} \rangle^2 + Z_t^2) \leq \phi(m_t)(1 - m_t^2)^2 + O\left(\frac{d}{n_b}\right).\] 
\end{proof}


\paragraph{Putting it all together} We now combine all the previous bounds into a unique proposition.
\begin{proposition}\label{2406:prop:app:diff_ineq_estimates}
    Let $T \geq 0$. There exists constants $c, C > 0$ such that
    \begin{equation}
        \begin{aligned}
            &\mathbb{P}\left(m_t \geq m_0 + \sum_{s=0}^{t-1} \Phi_{\mathrm{drift}}(m_s) - C\Phi_{\mathrm{noise}}(m_s) - C K(T) \quad \forall t \leq T \quad \big\vert\quad \bigcap_{t\leq T} \cE_t \right) \\
            &\geq 1 - Te^{-c\log(n)^2}
        \end{aligned}
    \end{equation}
    where $\Phi_{\mathrm{drift}}$ and $\Phi_{\mathrm{noise}}$ are given by
    \begin{align}
        \Phi_{\mathrm{drift}}(m) &= \gamma (1-m^2) \phi(m), \\
        \Phi_{\mathrm{noise}}(m) &=  \gamma^2 m(1 - m^2)\phi(m)^2 + \gamma^2 m \frac{d \log(d)^C}{n_b} +  \gamma^3 (1 - m^2)^{3/2}\phi(m)^3,  \\
        K(T) &= \frac{ \gamma \log(d)^C \sqrt{T}}{\sqrt{n_b}} +  \gamma^3 T\frac{d \log(d)^C}{n_b}.
    \end{align}
\end{proposition}

\begin{proof}
    By summing the inequality of Lemma \ref{2406:lem:app:difference_m_determ} for $0 \leq s \leq t-1$, we get
    \begin{equation}
        m_{t+1} \geq m_0 - \sum_{s=0}^{t-1} \gamma c_s \langle \vec{w}^\star,\vec{g}_s^\bot \rangle - \frac{\gamma^2 c_s^2 m_s \norm{\vec{g}_s^\bot}^2}{2} - \frac12\gamma^3 c_s^3 \norm{\vec{g}_s^\bot}^3.
    \end{equation}
    The linear term is handled using the martingale decomposition \eqref{2406:eq:app:martingale_decomp} combined with the expectation computation of \eqref{2406:eq:app:drift_term} and the bound of Lemma \ref{2406:lem:app:bound_martingale} on the martingale contribution. The terms in $\gamma^2$ and $\gamma^3$ follow from Lemma \ref{2406:lem:app:bound_corrections}, as well as Lemma \ref{2406:lem:app:bound_martingale} with $T = 1$, which implies that
    \[ |\langle \vec{w}^\star, \vec{g}_t^\bot \rangle \leq C \log(n)^C \]
    for some $C > 0$.
    Finally, under the events $\cE_s$ for $s\leq t$, we can replace every occurence of $c_s$ by either $1/2$ or $2$ depending on the sign of the corresponding term.
\end{proof}





\paragraph{Proof of Proposition \ref{2406:prop:app:diff_ineq_m}} The expressions of Lemma \ref{2406:lem:app:exp_g} imply the following expansions near 0:
\begin{align} 
\phi(m) &= c_{\ell}c_{\ell}^\star m^{\ell-1} + O(m^\ell) & \psi^{\mathrm{corr}}(m) &=  O(m^\ell) & \psi^{\mathrm{sq}}(m) = - c^{\mathrm{sq}} + O(m^\ell) \label{2406:eq:app:phi_psi_bounds}
\end{align} 
As a result, there exists an $\eta > 0$ and constants $C, c > 0$ such that for any $m \leq \eta$,
\begin{align*}
    c m^{\ell - 1} &\leq \phi(m) \leq C m^{\ell - 1} & |\psi^{\mathrm{corr}}(m)| &\leq C m^{\ell} & |\psi^{\mathrm{sq}}(m)| \leq C.
\end{align*}
We first lower bound the drift inequality of Proposition \ref{2406:prop:app:diff_ineq_estimates}. Whenever $m_t \geq \eta$, we have
\[ \Phi_{\mathrm{drift}}(m_t) \geq p(\gamma m_t^{\ell - 1}) - C\gamma^2\frac{d}{n_b},  \]
where $p(x) = x - C(x^2 + x^3)$. Define $\varepsilon > 0$ such that $p(x) > x/2$ on $[0, \varepsilon]$, so when $t \leq t_{\eta}^+ \wedge \tilde t_{\gamma,\varepsilon}^+$
\[ \Phi_{\mathrm{drift}}(m_t) \geq c\gamma m_t^{\ell-1} - C \gamma^2 \frac{d}{n_b} m.\]
When $\gamma \leq c_\gamma n_b d^{-\ell/2} \log(d)^{-C_\gamma}$,
\[ \gamma \frac{d \log(d)^C }{n_b} m \leq c_\gamma (\sqrt{d})^{\ell-2} \leq \frac{cm^{\ell-1}}2 \]
when $t \leq t_{\kappa/2\sqrt{d}}^-$, $C_\gamma \geq C$ and $c_\gamma \leq c/2 (\kappa/2)^{\ell-2}$.

Having shown $\Phi_{\mathrm{drift}}(m_t) \geq c m_t^{\ell-1}$, it remains to handle the constant terms in Proposition \ref{2406:prop:app:diff_ineq_estimates}. We can compute directly $K(t_{\max})$, which yields
\[K(t_{\max}) = \frac{\log(d)^{C - C_{\max}/2}}{C_{\max}\sqrt{d}} + C\gamma\log(d)^{C - C_{\max}}\frac d{n_b} \leq \frac{\log(d)^{C - C_{\max}/2}}{C_{\max}\sqrt{d}} + \frac{Cc_\gamma}{\log(d)^{C_{\max}}} d^{-\frac{\ell}{2}} \leq \frac{m_0}{4}, \]
by choosing $c_\gamma$ small enough and $C_{\max}$ large enough.

Finally, we need to show that the events $\cE_t$ occur with high probability. This is covered by the following lemma:
\begin{lemma}
    Let $T \geq 0$. The following bounds hold:
    \begin{itemize}
        \item for the square loss, if $\gamma \leq c_\gamma$ for small enough $c_\gamma$,
        \[ \mathbb{P}\left(\cE_t \text{ holds for all } t \leq t_\eta^+ \wedge T \right) \geq 1 - Te^{-c\log(d)^2}; \]
        \item for the correlation loss, for small enough $\varepsilon$,
        \[ \mathbb{P}\left(\cE_t \text{ holds for all } t \leq t_\eta^+ \wedge \tilde t_{\gamma, \varepsilon}^+ \wedge T \right) \geq 1 - Te^{-c\log(d)^2}. \]
    \end{itemize}
\end{lemma}

\begin{proof}
    We begin with the case of the square loss. From the expression of the gradient expectation in \eqref{2406:eq:app:exp_g}, and the estimates \eqref{2406:eq:app:phi_psi_bounds},
    \[ \Et{\langle \vec{w}_t, \vec{g}_t \rangle} = \psi^{\mathrm{sq}}(m_t) - \phi(m_t)m_t = O(1),\]
    whenever $t \leq t_\eta^+$. Lemma \ref{2406:lem:app:orlicz_concentration} applied to $\langle \vec{w}_t, \vec{g}_t \rangle$ implies that with probability $1 - ce^{-c\log(n)^2}$
    \[ |\langle \vec{w}_t, \vec{g}_t \rangle| \leq |\Et{\langle \vec{w}_t, \vec{g}_t \rangle}| +\frac{C\log(d)^C}{\sqrt{n_b}} = O(1) \]
    whenever $\mu > 0$. As a result, if $\gamma \leq c_\gamma$ for $c_\gamma$ small enough, $\cE_t$ holds. The proof for the correlation loss proceeds identically, noting this time that
    \[ \Et{\langle \vec{w}_t, \vec{g}_t \rangle} = \psi^{\mathrm{corr}}(m_t) - \phi(m_t)m_t = O(1) \]
    whenever $t \leq t_\eta^+ \wedge \tilde t_{\gamma, \varepsilon}^+$.
\end{proof}

\subsection{From the linear regime to one-step recovery}
We now prove Theorems \ref{2406:thm:main:sgd_weak_recovery} and \ref{2406:thm:main:no_yhat_weak_recovery}. We focus on the case of the correlation loss; the square loss is identical apart from the additional $\gamma \leq c_\gamma$ requirement of Proposition \ref{2406:prop:app:diff_ineq_m}. 

We first assume that $n_b = O(d^{\ell-1})$ and
\begin{equation}\label{2406:eq:app:bound_gamma}
    \gamma \leq  c_\gamma \log(d)^{-C_\gamma} \min\left( n_b d^{-\left(\frac\ell2 \vee 1 \right)}, \sqrt{\frac{n_b}d} \right)
\end{equation} 
for constants $c_\gamma, C_\gamma$ to be chosen later. In particular, the condition $\gamma < c_\gamma n_b d^{-\ell/2}$ of Proposition \ref{2406:prop:app:diff_ineq_m} is satisfied.

\paragraph{The linear regime} The first part of the proof proceeds as in \cite{arous2021online}. Define the function
\begin{equation}
    t_\ell(d) = \begin{cases}
        1 &\text{if } \ell = 1 \\
        \log(d)  &\text{if } \ell = 2 \\
        d^{\frac{\ell}2-1}  &\text{if } \ell > 2
    \end{cases}, 
\end{equation}
and $t_{\mathrm{conv}} = \max(1, \gamma^{-1}t_\ell(d))$.
Proposition \ref{2406:prop:app:diff_ineq_m} as well as Section 5 from \cite{arous2021online} implies the following lemma:
\begin{lemma} \label{2406:lem:app:boundongamma}
    There exists a constant $C > 0$ such that if $t_{\max} \geq C t_{\mathrm{conv}}$, then with probability at least $1 - ce^{-c\log(n)^2}$,
    \begin{equation}
        \tilde t_{\gamma, \varepsilon} \wedge t_\eta^+ \leq Ct_{\mathrm{conv}}.
    \end{equation}
\end{lemma}
We therefore only need to check the condition $t_{\max} \geq C t_{\mathrm{conv}}$.
Plugging the expression for $\gamma$ and $t_{\max}$, we get
\begin{align*}
    \frac{\gamma t_{\max}}{t_\ell(d)} &\geq \frac{n_b}{\gamma C_{\max} \log(d)^{C_{\max}+1} \gamma d^{1 + (\ell/2 - 1)\vee 0}} \geq \frac1{c_\gamma C_{\max}} \log(d)^{C_\gamma - C_{\max} - 1} \\
     t_{\max} &\geq \frac{n_b}{C_{\max}\log(d)^{C_{\max}} d \gamma^2} \geq \frac1{c_\gamma C_{\max}} \log(d)^{2C_\gamma - C_{\max} - 1}
\end{align*} 
Whenever $n_b = O(d^{\ell-1})$, by decreasing $c_\gamma$ and increasing $C_\gamma$, for large enough $d$ and any constant $C$ we have
\[\frac{\gamma t_{\max}}{t_\ell(d)} \leq C \quad \text{and} \quad t_{\max} \geq C, \]
as requested in Lemma \ref{2406:lem:app:boundongamma}.

Whenever $\gamma \eta^{\ell-1} \leq \varepsilon$, we have $\tilde t_{\gamma, \varepsilon} \geq t_\eta^+$ and hence the proof of Theorem \ref{2406:thm:main:no_yhat_weak_recovery} is complete. It thus remains to treat the converse case. Note that the latter only happens in the correlation loss case, since we can always choose $c_\gamma$ such that $c_\gamma \eta^{\ell-1} \leq \varepsilon$; the dynamics of square loss SGD are therefore only in the linearized regime.

\paragraph{One-step recovery above $\tilde t_{\gamma, \varepsilon}$} We now treat the case where $\gamma \eta^{\ell-1} \geq \varepsilon$. For simplicity, let $t = \tilde t_{\gamma, \varepsilon}$, then Lemmas \ref{2406:lem:app:bound_corrections} and \ref{2406:lem:app:bound_martingale} imply that with probability $1 - ce^{-c\log(n)^2}$
\begin{align*}
    \norm{\vec{g}_{t}}^2 &\leq C\left((|\phi(m_t)| + |\psi^{\mathrm{corr}}(m_t)|)^2 + \frac{d}{n_b}\right) \leq C \left(m_t^{2\ell-2} + \frac{d}{n_b}\right) \\
    \langle \vec{g}_t, \vec{w}^\star \rangle &= \phi(m_t) + m_t\psi^{\mathrm{corr}}(m_t) + O\left( \frac{\polylog(d)}{\sqrt{n_b}} \right) \geq cm_t^{\ell - 1} + O\left( \frac{\polylog(d)}{\sqrt{n_b}} \right)
\end{align*}
By definition of $t$, we have $1 \leq \varepsilon^{-1}\gamma m_t^{\ell-1}$, and whenever $\gamma \leq c_\gamma\sqrt{n_b/d}$ one has
\[ \frac\gamma{\sqrt{n_b}} \leq c_\gamma d^{-1/2} \quad \text{and} \quad \gamma^2 \frac{d}{n_b} \leq c_\gamma^2.\]
Therefore, for large enough $d$,
\begin{align*} 
m_t + \gamma \langle \vec{g}_t, \vec{w}^\star \rangle &\geq  \gamma \langle \vec{g}_t, \vec{w}^\star \rangle \geq c\gamma m_t^{\ell-1}\\
\norm{\vec{w}_t + \gamma \vec{g}_t} &\leq 1 + \gamma \norm{\vec{g}_t} \leq 1 + C \gamma \left(m_t^{\ell-1} + \frac{d}{n_b}\right) \leq (C + \varepsilon^{-1} + c_\gamma^2\varepsilon^{-1}) \gamma m_t^{\ell-1} 
\end{align*}
But by taking the scalar product of Equation \eqref{2406:eq:main:gd_update_weights} with $\vec{w}^\star$,
\[ m_{t+1} = \frac{m_t + \gamma \langle \vec{g}_t, \vec{w}^\star \rangle}{\norm{\vec{w}_t + \gamma \vec{g}_t}} \geq \frac{c}{C + \varepsilon^{-1} + c_\gamma^2\varepsilon^{-1}} =: \eta'. \]
Theorem \ref{2406:thm:main:no_yhat_weak_recovery} ensues by redefining $\eta = \min(\eta, \eta').$
